\textbf{1.} Check reflexivity, symmetry, and transitivity using nonzero scalars.

\textbf{2.} A standard chart is present exactly when its corresponding homogeneous coordinate is nonzero.

\textbf{3.} Normalize the first homogeneous coordinate to one.

\textbf{4.} Use \(x_0/x_1=1/u\) and \(x_2/x_1=v/u\).

\textbf{5.} Both coordinates are nonzero on the overlap, so localize at the affine coordinate.

\textbf{6.} Only the homogeneous prime \((0)\) survives and its degree-zero chart ring is \(k\).

\textbf{7.} Set \(x_0=0\) and retain the remaining homogeneous coordinates.

\textbf{8.} Use one nonempty affine chart for the lower bound and the affine cover for the upper bound.

\textbf{9.} Only \(x_3,\ldots,x_n\) remain.

\textbf{10.} Projectivize the span of the two coordinate vectors.

\textbf{11.} Use \(\mathbb R[t]/(t^2+1)\).

\textbf{12.} Base change \(k[t]\) and \(k[u]\), then check \(u=t^{-1}\) is unchanged.

\textbf{13.} Use \(\mathcal O(1)=\widetilde{S(1)}\) and localize at \(x_i\).

\textbf{14.} Compare \(x_i^d\) and \(x_j^d\) in the common homogeneous localization.

\textbf{15.} Degree-one homogeneous polynomials are linear forms.

\textbf{16.} Count monomials of degree two in \(n+1\) variables.

\textbf{17.} Set the chosen chart coordinate equal to one and divide by the appropriate power.

\textbf{18.} Introduce \(x_0\) and replace \(x,y\) by \(x_1/x_0,x_2/x_0\) before clearing denominators.

\textbf{19.} Every degree-\(d\) monomial scales by \(\lambda^d\).

\textbf{20.} Substitute \(X=s^2,Y=st,Z=t^2\).

\textbf{21.} Track scalars \(\lambda\) and \(\mu\) in the two factors.

\textbf{22.} Substitute \(z_{00}=su,z_{01}=sv,z_{10}=tu,z_{11}=tv\).

\textbf{23.} Take \(d=0\).

\textbf{24.} The affine recovery theorem applies only to affine schemes.